\documentclass[12pt,a4paper]{report}

% =========================
% Packages
% =========================
\usepackage[utf8]{inputenc}
\usepackage[T1]{fontenc}
\usepackage[french]{babel}
\usepackage{geometry}
\usepackage{graphicx}
\usepackage{fancyhdr}
\usepackage{titlesec}
\usepackage{xcolor}
\usepackage{float}
\usepackage{enumitem}
\usepackage[most]{tcolorbox}
\usepackage{hyperref}

% =========================
% Geometry
% =========================
\geometry{
    left=2.5cm,
    right=2.5cm,
    top=2.5cm,
    bottom=2.5cm
}

% =========================
% Colors
% =========================
\definecolor{ensablue}{RGB}{0,51,102}

% =========================
% Hyperref setup (après les autres packages)
% =========================
\hypersetup{
    colorlinks=true,
    linkcolor=ensablue,
    filecolor=magenta,
    urlcolor=blue
}

% =========================
% Header and footer
% =========================
\pagestyle{fancy}
\fancyhf{}
\fancyhead[L]{\small\textit{Application de Réservation de Billets}}
\fancyhead[R]{\thepage}
\renewcommand{\headrulewidth}{0.5pt}
\setlength{\headheight}{15pt}

% =========================
% Chapter and section formatting
% =========================
\titleformat{\chapter}[display]
{\normalfont\Large\bfseries\color{ensablue}}
{\chaptertitlename\ \thechapter}{10pt}{\Large}
\titlespacing*{\chapter}{0pt}{-20pt}{20pt}

\titleformat{\section}
{\normalfont\large\bfseries\color{ensablue}}
{\thesection}{1em}{}

\titleformat{\subsection}
{\normalfont\normalsize\bfseries\color{ensablue}}
{\thesubsection}{1em}{}

\begin{document}

% ============================================
% PAGE DE GARDE
% ============================================
\begin{titlepage}
    \begin{center}
        \vspace*{0.5cm}
        \includegraphics[width=0.4\textwidth]{image.png}

        \vspace{1cm}

        {\large \textbf{Université Abdelmalek Essaâdi}}\\
        {\large \textbf{École Nationale des Sciences Appliquées de Tétouan}}\\

        \vspace{0.5cm}
        \rule{\textwidth}{1pt}
        \vspace{0.8cm}

        {\color{ensablue}\textbf{FILIÈRE : GÉNIE INFORMATIQUE 2}}\\[0.3cm]
        {\large Module : Microservices}\\

        \vspace{1.2cm}

        \begin{tcolorbox}[
            colback=gray!10,
            colframe=ensablue,
            width=0.9\textwidth,
            arc=8mm,
            boxrule=0.8pt,
            drop shadow
        ]
        \centering
        \Large\bfseries
        Rapport d'avancement\\
        Microservice ``Coordination des transports'' (G4)
        \end{tcolorbox}

        \vspace{1.5cm}

        \begin{minipage}[t]{0.45\textwidth}
            \textbf{Réalisé par :}\\[0.3cm]
            EL ASSAL DOUAE\\
            EL MAKOUDI YOUNES\\
            MAGNE TSAFACK LYDIVINE MERVEILLE\\
        \end{minipage}
        \hfill
        \begin{minipage}[t]{0.45\textwidth}
            \raggedleft
            \textbf{Encadré par :}\\[0.3cm]
            Pr. BESRI\\
            \vspace{0.5cm}
        \end{minipage}

        \vfill

        \rule{\textwidth}{1pt}\\
        \vspace{0.3cm}
        {\large \textbf{Année Universitaire : 2025--2026}}

    \end{center}
\end{titlepage}

% ============================================
% TABLE DES MATIÈRES
% ============================================
\tableofcontents
\newpage

% ============================================
% CHAPITRE 1
% ============================================
\chapter{Présentation du microservice}

\section{Description du sous-système}
Le microservice \textbf{G4 -- Coordination des Transports} est responsable de la
coordination opérationnelle du réseau SGITU. Il centralise la gestion des entités
métier nécessaires à l'exploitation quotidienne : lignes, trajets, arrêts,
horaires, affectations véhicule/ligne, missions et événements de coordination
(retard, déviation, panne, incident, annulation).

Son rôle est de fournir une vue cohérente de l'état opérationnel et de permettre
le pilotage des missions de transport en temps réel.

\section{Positionnement dans SGITU}
Dans l'architecture microservices SGITU, G4 occupe une position centrale entre
les services de planification/exploitation et les services d'intégration temps réel.

\begin{itemize}
  \item \textbf{Exposition API} : endpoints REST exposés via la Gateway G10.
  \item \textbf{Interaction suivi des véhicules G7} : consommation des positions véhicules pour le suivi opérationnel et la détection de situations anormales.
  \item \textbf{Interaction gestion des incidents G9} : consommation des incidents transport et transformation en événements de coordination.
  \item \textbf{Interaction Billetterie G1} : publication du cycle de vie des missions (Kafka) selon le contrat inter-services.
  \item \textbf{Interaction Notifications} : relais d'alertes vers le service de notification.
\end{itemize}

Ainsi, G4 agit comme couche de \textit{coordination métier} entre les données de
référence, les flux temps réel et les actions opérationnelles.

\section{Fonctionnalités métier implémentées}
Les fonctionnalités implémentées à ce stade sont les suivantes :

\begin{enumerate}
  \item \textbf{Gestion du référentiel transport}
  \begin{itemize}
    \item CRUD des lignes ;
    \item CRUD des trajets et consultation des arrêts par trajet ;
    \item CRUD des arrêts et filtrage par ligne ;
    \item CRUD des horaires.
  \end{itemize}

  \item \textbf{Gestion opérationnelle}
  \begin{itemize}
    \item gestion des affectations véhicule/ligne ;
    \item planification et suivi des missions ;
    \item transitions de mission (en cours, clôture, annulation) .
  \end{itemize}

  \item \textbf{Gestion des événements de coordination}
  \begin{itemize}
    \item enregistrement d'événements métier ;
    \item détection de retard, déviation, panne et incident ;
    \item annulation de mission sur décision opérationnelle.
  \end{itemize}

  \item \textbf{Intégration et supervision}
  \begin{itemize}
    \item publication Kafka des événements de cycle de vie mission ;
    \item endpoints de supervision (\texttt{/api/g4/health}, statut opérateur, logs) ;
    \item documentation OpenAPI/Swagger et collection Postman de tests.
  \end{itemize}
\end{enumerate}

% ============================================
% CHAPITRE 2
% ============================================
\chapter{Conception}

\section{Identification des endpoints REST}
Les endpoints du microservice G4 sont exposés principalement sous le préfixe \texttt{/api/g4/**}.  
Le service de notification est exposé via \texttt{/api/notifications/send}.  
Les endpoints de supervision et d'état opérateur sont également disponibles pour le monitoring du système.

\section{Choix des méthodes HTTP}
Le design REST suit les conventions standards :
\begin{itemize}
    \item \texttt{GET} : consultation des ressources ;
    \item \texttt{POST} : création de ressource ou action métier ;
    \item \texttt{PUT} : mise à jour d'une ressource ;
    \item \texttt{DELETE} : suppression d'une ressource.
\end{itemize}

\section{Formats d'échange}
Le format d'échange principal est \textbf{JSON}.  
Les contrats d'entrée/sortie sont portés par des DTOs avec validations.  
L'API est documentée via \textbf{Swagger/OpenAPI}.

\section{Gestion des erreurs}
Les codes HTTP appliqués sont :
\begin{itemize}
    \item \texttt{200} / \texttt{201} : succès ;
    \item \texttt{400} : requête invalide (validation) ;
    \item \texttt{401} : non authentifié ;
    \item \texttt{403} : interdit (rôle insuffisant) ;
    \item \texttt{404} : ressource introuvable ;
    \item \texttt{500} : erreur interne.
\end{itemize}

% ============================================
% CHAPITRE 3
% ============================================
\chapter{Implémentation}

\section{Architecture en couches}
Le microservice suit une architecture en couches :
\begin{itemize}
    \item \texttt{controller} : exposition REST ;
    \item \texttt{service} : logique métier ;
    \item \texttt{repository} : accès aux données via JPA ;
    \item \texttt{entity} : modèle de persistance ;
    \item \texttt{dto/mapper} : transformation entrée/sortie ;
    \item \texttt{config/security} : sécurité JWT, configuration Kafka et technique.
\end{itemize}

\section{Technologies utilisées}
\begin{itemize}
    \item Java, Spring Boot ;
    \item Spring Web, Spring Data JPA, Spring Security ;
    \item JWT (authentification locale de développement) ;
    \item Kafka (publication/consommation d'événements) ;
    \item PostgreSQL (runtime) ;
    \item Swagger/OpenAPI ;
    \item Docker, Docker Compose.
\end{itemize}
\section{Organisation du projet}
Le projet suit une structure Maven standard avec une séparation claire des responsabilités :
\begin{itemize}
    \item \texttt{controller} : exposition des endpoints REST ;
    \item \texttt{service} : implémentation de la logique métier ;
    \item \texttt{repository} : accès aux données via Spring Data JPA ;
    \item \texttt{entity} : modèle de persistance ;
    \item \texttt{dto} : contrats d'échange API (request/response) ;
    \item \texttt{mapper} : transformation entre entités et DTO ;
    \item \texttt{integration} : échanges avec les services externes (G1, G5, G7, G9) ;
    \item \texttt{config} / \texttt{security} : configuration technique, JWT et sécurité ;
    \item \texttt{exception} : gestion centralisée des erreurs.
\end{itemize}

% ============================================
% CHAPITRE 4
% ============================================
\chapter{Documentation API (Swagger / OpenAPI)}

\section{Documentation complète}
La documentation Swagger UI couvre les endpoints d'authentification, de gestion métier (CRUD + actions), de supervision et d'intégration.

\section{Requêtes et réponses}
Chaque endpoint dispose de payloads JSON structurés :
\begin{itemize}
    \item Request DTO avec validations ;
    \item Response DTO cohérente avec le domaine.
\end{itemize}

\section{Cas d'utilisation couverts}
\begin{itemize}
    \item Création et gestion de lignes, trajets, arrêts, horaires ;
    \item Affectation opérationnelle véhicule/ligne ;
    \item Création et suivi des missions ;
    \item Détection/déclaration d'événements de coordination ;
    \item Clôture et annulation de mission.
\end{itemize}

\section{Accès OpenAPI}
\begin{itemize}
    \item \texttt{/swagger-ui.html}
    \item \texttt{/v3/api-docs}
\end{itemize}

\begin{figure}[H]
  \centering
  \includegraphics[width=0.92\textwidth]{fig_4_1_swagger_overview.png}
  \caption{Vue d'ensemble Swagger UI montrant les endpoints exposés du microservice G4.}
  \label{fig:fig_4_1}
\end{figure}

\begin{figure}[H]
  \centering
  \includegraphics[width=0.92\textwidth]{fig_4_2_swagger_request_response.png}
  \caption{Exécution d'un endpoint dans Swagger avec payload de requête JSON et réponse HTTP.}
  \label{fig:fig_4_2}
\end{figure}

\begin{figure}[H]
  \centering
  \includegraphics[width=0.92\textwidth]{fig_4_3_openapi_json.png}
  \caption{Extrait du contrat OpenAPI (\texttt{/v3/api-docs}) montrant version, paths et sécurité JWT (\texttt{bearerAuth}).}
  \label{fig:fig_4_3}
\end{figure}

% ============================================
% CHAPITRE 5
% ============================================
\chapter{Conteneurisation (Docker)}

\section{Introduction}
Afin d'assurer la portabilité du microservice et la reproductibilité de son exécution,
la conteneurisation repose sur une image Docker dédiée à G4. Les dépendances
d'infrastructure (PostgreSQL et Kafka) sont orchestrées séparément via Docker Compose.

\section{Dockerfile}
Le Dockerfile est de type \textit{multi-stage}, ce qui signifie qu'il contient deux étapes :
\begin{itemize}
    \item une étape de build (compilation et packaging Maven) ;
    \item une étape runtime (exécution du JAR dans une image plus légère).
\end{itemize}
Cette approche réduit la taille de l'image finale et isole les dépendances de compilation de l'environnement d'exécution.

\begin{verbatim}
# Build
FROM eclipse-temurin:17-jdk-alpine AS build
WORKDIR /app
COPY .mvn .mvn
COPY mvnw pom.xml ./
COPY src src
RUN chmod +x mvnw && ./mvnw -q -DskipTests package

# Run
FROM eclipse-temurin:17-jre-alpine
WORKDIR /app
RUN addgroup -S sgitu && adduser -S g4 -G sgitu
COPY --from=build /app/target/service-coordination-transport-0.0.1-SNAPSHOT.jar app.jar
RUN chown g4:sgitu app.jar
USER g4
EXPOSE 8084
ENTRYPOINT ["sh", "-c", "java $JAVA_OPTS -jar /app/app.jar"]
\end{verbatim}

\section{Image Docker}
L'image Docker du service est générée localement avec succès à partir du Dockerfile multi-stage.

\begin{figure}[H]
  \centering
  \includegraphics[width=0.92\textwidth]{fig_5_1_docker_build_success.png}
  \caption{Build Docker réussi du microservice G4 (image générée et taguée localement).}
  \label{fig:fig_5_1}
\end{figure}
\section{Exécution}
Le service est exécutable via Docker Compose et exposé sur le port \texttt{8084}.  
La validation d'exécution est démontrée par :
\begin{itemize}
    \item le statut \texttt{Running} des conteneurs dans Docker Compose ;
    \item l'accès à l'interface Swagger sur \texttt{http://localhost:8084/swagger-ui.html}.
\end{itemize}

% ============================================
% CHAPITRE 6
% ============================================
\chapter{Orchestration (docker-compose)}

\section{Introduction}
L'orchestration locale est unifiée dans un seul fichier \texttt{docker-compose.yml},
ce qui permet de démarrer rapidement un environnement cohérent de test et d'intégration.

\section{docker-compose.yml (contenu exact)}
\begin{verbatim}
services:
  kafka:
    image: apache/kafka:3.7.0
    container_name: sgitu-kafka
    ports:
      - "9092:9092"
    environment:
      KAFKA_NODE_ID: 1
      KAFKA_PROCESS_ROLES: broker,controller
      KAFKA_CONTROLLER_QUORUM_VOTERS: 1@localhost:9093
      KAFKA_LISTENERS: PLAINTEXT://:9092,CONTROLLER://:9093
      KAFKA_ADVERTISED_LISTENERS: PLAINTEXT://localhost:9092
      KAFKA_LISTENER_SECURITY_PROTOCOL_MAP: CONTROLLER:PLAINTEXT,PLAINTEXT:PLAINTEXT
      KAFKA_CONTROLLER_LISTENER_NAMES: CONTROLLER
      KAFKA_OFFSETS_TOPIC_REPLICATION_FACTOR: 1
      KAFKA_TRANSACTION_STATE_LOG_REPLICATION_FACTOR: 1
      KAFKA_TRANSACTION_STATE_LOG_MIN_ISR: 1
      KAFKA_AUTO_CREATE_TOPICS_ENABLE: "true"
    healthcheck:
      test: ["CMD-SHELL", "bash -c '</dev/tcp/127.0.0.1/9092'"]
      interval: 10s
      timeout: 5s
      retries: 12
    networks:
      - sgitu-net

  postgres:
    image: postgres:16-alpine
    container_name: sgitu-g4-postgres
    environment:
      POSTGRES_DB: sgitu_g4
      POSTGRES_USER: g4
      POSTGRES_PASSWORD: g4
    ports:
      - "5434:5432"
    volumes:
      - sgitu_g4_pgdata:/var/lib/postgresql/data
    healthcheck:
      test: ["CMD-SHELL", "pg_isready -U g4 -d sgitu_g4"]
      interval: 5s
      timeout: 5s
      retries: 10
    networks:
      - sgitu-net

  g4-coordination:
    build: .
    container_name: sgitu-g4-coordination
    ports:
      - "8084:8084"
    environment:
      SPRING_PROFILES_ACTIVE: docker
      SERVER_PORT: "8084"
      JPA_DDL_AUTO: update
      SGITU_KAFKA_ENABLED: "true"
      SPRING_KAFKA_BOOTSTRAP_SERVERS: "kafka:9092"
      SGITU_G9_INCIDENT_TOPIC: "${SGITU_G9_INCIDENT_TOPIC:-incident.transport.topic}"
      SGITU_G7_POSITIONS_TOPIC: "${SGITU_G7_POSITIONS_TOPIC:-vehicule-positions}"
      SGITU_G1_MISSION_TOPIC: "${SGITU_G1_MISSION_TOPIC:-missions-lifecycle}"
      SGITU_G9_CONSUMER_GROUP: "${SGITU_G9_CONSUMER_GROUP:-g4-coordination-g9}"
      SGITU_G7_CONSUMER_GROUP: "${SGITU_G7_CONSUMER_GROUP:-g4-coordination-g7}"
      SGITU_G10_URL: "${SGITU_G10_URL:-http://g10-gateway:8080}"
      SGITU_G5_URL: "${SGITU_G5_URL:-http://g5-notifications:8085}"
      SGITU_G1_URL: "${SGITU_G1_URL:-http://g1-billetterie:8081}"
      SGITU_G7_URL: "${SGITU_G7_URL:-http://g7-tracking:8087}"
      SGITU_G9_URL: "${SGITU_G9_URL:-http://g9-incidents:8089}"
      SGITU_INTEGRATION_G5_NOTIFICATION_PATH: "${SGITU_INTEGRATION_G5_NOTIFICATION_PATH:-/api/notifications/send}"
      SGITU_G1_MISSION_LIFECYCLE_PATH: "${SGITU_G1_MISSION_LIFECYCLE_PATH:-/api/internal/missions/events}"
    depends_on:
      postgres:
        condition: service_healthy
      kafka:
        condition: service_healthy
    networks:
      - sgitu-net

volumes:
  sgitu_g4_pgdata:

networks:
  sgitu-net:
    name: sgitu-network
    driver: bridge
\end{verbatim}

\section{Interaction inter-microservices}
Le compose prépare les variables d'environnement nécessaires à l'intégration avec la Gateway (G10), les notifications (G5), la billetterie (G1), le suivi véhicule (G7) et les incidents (G9).

\section{Simulation distribuée}
L'environnement local simule un contexte distribué réaliste avec service applicatif, base PostgreSQL et broker Kafka.

\begin{figure}[H]
  \centering
  \includegraphics[width=0.92\textwidth]{fig_6_1_compose_running_3_3.png}
  \caption{Environnement Docker Compose G4 démarré avec succès : \texttt{sgitu-kafka}, \texttt{sgitu-g4-postgres} et \texttt{sgitu-g4-coordination} en statut \textit{Running (3/3)}.}
  \label{fig:fig_6_1}
\end{figure}

\begin{figure}[H]
  \centering
  \includegraphics[width=0.92\textwidth]{fig_6_2_docker_desktop.png}
  \caption{Docker Desktop : projet \texttt{service-coordination-transport} avec les conteneurs \texttt{sgitu-kafka} (\texttt{apache/kafka:3.7.0}, port \texttt{9092}), \texttt{sgitu-g4-coordination} (port \texttt{8084}) et \texttt{sgitu-g4-postgres} (\texttt{postgres:16-alpine}, port hôte \texttt{5434}) en cours d'exécution.}
  \label{fig:fig_6_2}
\end{figure}

% ============================================
% CHAPITRE 7
% ============================================
\chapter{Sécurité (JWT, Spring Security)}

\section{Authentification (JWT)}
Endpoint de login local : \texttt{POST /api/auth/login}, retour d'un token \texttt{Bearer JWT}.
\section{Protection des endpoints}
Les endpoints publics sont limités à :
\begin{itemize}
    \item \texttt{/api/auth/login} ;
    \item \texttt{/api/g4/health} ;
    \item \texttt{/swagger-ui.html} et \texttt{/v3/api-docs}.
\end{itemize}
Tous les autres endpoints sont protégés par Spring Security + JWT.

\section{Gestion des accès (rôles)}
Le contrôle d'accès est défini autour de trois rôles métier avec séparation fonctionnelle.

\subsection*{ROLE\_G4\_OPERATOR}
Rôle de consultation et de supervision :
\begin{itemize}
  \item lecture des ressources métier (\texttt{GET /api/g4/**}) ;
  \item consultation du statut opérateur (\texttt{GET /api/v1/operator/status}) ;
  \item consultation des logs de supervision (\texttt{GET /api/g4/logs}).
\end{itemize}
Ce rôle ne peut ni créer, ni modifier, ni supprimer des ressources.

\subsection*{ROLE\_G4\_DISPATCHER}
Rôle de pilotage opérationnel :
\begin{itemize}
  \item consultation des ressources métier ;
  \item gestion opérationnelle des affectations, missions et événements de coordination ;
  \item actions métier (création, mise à jour, annulation/clôture selon endpoints opérationnels) ;
  \item envoi de notifications via \texttt{POST /api/notifications/send}.
\end{itemize}
Ce rôle n'est pas destiné à administrer le référentiel structurel du réseau
(lignes, trajets, arrêts, horaires).

\subsection*{ROLE\_G4\_ADMIN}
Rôle d'administration globale :
\begin{itemize}
  \item accès complet à l'ensemble des fonctionnalités du microservice G4 ;
  \item gestion du référentiel structurel (lignes, trajets, arrêts, horaires) ;
  \item supervision globale et gouvernance des opérations.
\end{itemize}

\subsection*{Règles d'accès fonctionnelles}
\begin{itemize}
  \item \textbf{Consultation (\texttt{GET})} : \texttt{ROLE\_G4\_OPERATOR}, \texttt{ROLE\_G4\_DISPATCHER}, \texttt{ROLE\_G4\_ADMIN}.
  \item \textbf{Référentiel structurel} (\texttt{POST}/\texttt{PUT}/\texttt{DELETE} sur lignes, trajets, arrêts, horaires) : \texttt{ROLE\_G4\_ADMIN}.
  \item \textbf{Opérationnel} (\texttt{POST}/\texttt{PUT}/\texttt{DELETE} sur affectations, missions, événements) : \texttt{ROLE\_G4\_DISPATCHER}, \texttt{ROLE\_G4\_ADMIN}.
  \item \textbf{Notifications} (\texttt{POST /api/notifications/send}) : \texttt{ROLE\_G4\_DISPATCHER}, \texttt{ROLE\_G4\_ADMIN}.
\end{itemize}

\section{Supervision et monitoring}
Le microservice expose des endpoints dédiés au monitoring :
\begin{itemize}
  \item \texttt{GET /api/g4/health} : état de santé global du service ;
  \item \texttt{GET /api/v1/operator/status} : état opérationnel orienté exploitation ;
  \item \texttt{GET /api/g4/logs} : consultation des logs de supervision applicative.
\end{itemize}
Ces endpoints facilitent le suivi en temps réel, le diagnostic et l'observabilité du service en exploitation.

% ============================================
% CHAPITRE 8
% ============================================
\chapter{État d'avancement, difficultés, perspectives}

\section{Fonctionnalités finalisées}
\begin{itemize}
    \item Endpoints métiers principaux implémentés ;
    \item Persistance des données ;
    \item Auth JWT locale + sécurité par rôles ;
    \item Kafka producer/consumers ;
    \item Dockerfile + docker-compose unifié ;
    \item Swagger, Postman, tests Maven.
\end{itemize}

\section{Fonctionnalités en cours / à finaliser}
\begin{itemize}
    \item Alignement final de l'authentification inter-services avec le groupe Users/Gateway (source d'identité et format JWT commun) ;
    \item Validation finale des URLs et ports externes selon les conventions inter-groupes ;
    \item Harmonisation de la claim des rôles (\texttt{roles}/\texttt{authorities}) sur la plateforme.
\end{itemize}

\section{Difficultés rencontrées}
\begin{itemize}
    \item Alignement des rôles entre équipes ;
    \item Clarification de l'architecture JWT (Service Users vs Gateway) ;
    \item Image Kafka initiale indisponible, corrigée dans le compose unifié.
\end{itemize}

\section{Perspectives}
\begin{itemize}
    \item Finalisation de l'intégration sécurité globale SGITU ;
    \item Stabilisation des flux Kafka inter-groupes ;
    \item Renforcement des tests d'intégration de bout en bout.
\end{itemize}

\end{document}
